\documentclass[11pt]{article}
\usepackage[margin=1in]{geometry}
\usepackage[T1]{fontenc}
\usepackage[utf8]{inputenc}
\usepackage{lmodern}
\usepackage{microtype}
\usepackage{enumitem}
\usepackage{xcolor}
\usepackage{booktabs}
\usepackage{array}
\usepackage{longtable}
\usepackage{hyperref}
\usepackage{multirow} 
\usepackage{amsmath,amssymb}
\usepackage{graphicx}
\usepackage{rotating}
\setlength{\parindent}{0pt}
\setlength{\parskip}{0.6em}
\usepackage{xcolor}
\newcommand{\reviewer}[1]{\vspace{1em}{\Large\bfseries Reviewer #1}\par\vspace{0.5em}}
\newcommand{\question}[2]{\vspace{0.5em}\noindent\textbf{Question #1.} \emph{#2}\par}
\newcommand{\answer}{\noindent\textbf{Answer:}\par}
%\setlength{\parindent}{2em}   % Thụt đầu dòng
\setlength{\parskip}{1em}   % Khoảng cách giữa các đoạn

\begin{document}
\title{
    {\bfseries Response to the Editor and Reviewers' Comments: ID 6b81018f-265c-42a3-8426-972a6621e2e9}\\\vspace*{2cm}
    %
    {\Large  Adaptive Hybrid Evolutionary Search for Dynamic Pickup and Delivery Problems with Time Windows}\\\vspace*{2cm}
    %
   \vspace*{5cm} }

\author{}

\date{November 2026}
\maketitle\thispagestyle{empty}\newpage
\setlength{\parskip}{2ex plus0.5ex minus0.2ex} % only for cover letter

\thispagestyle{empty} \vspace{-10mm}

\noindent Dear Editor,

\medskip

\noindent
Please find enclosed the further revised version of our manuscript entitled
``Adaptive Hybrid Evolutionary Search for Dynamic Pickup and Delivery Problems with Time Windows''
(Manuscript ID: 6b81018f-265c-42a3-8426-972a6621e2e9).

\medskip

\noindent
We would like to sincerely thank you for your time and effort in handling our manuscript and coordinating the review process. We are also grateful to the reviewers for their careful evaluation and constructive comments.

\medskip

\noindent
Our responses to the comments are provided below.

\noindent The revised Related Work integrates the three studies suggested by Reviewer~4 and expands the discussion of methods relevant to the ICAPS benchmark.

\noindent Yours sincerely,

\noindent Tran Dinh Phuc, Nguyen Thi Tam, and Huynh Thi Thanh Binh\vspace{5mm}

\setlength{\parskip}{0pt} % remove excessive parskip again

\clearpage
\begin{center}
\textbf{\Large Authors' Response to Reviewer 1}
\end{center}

\begin{quote}\emph{I have no more comments.}\end{quote}

\answer
We thank Reviewer~1 for confirming that there are no further comments. We appreciate the reviewer's helpful feedback during the previous round.

\clearpage
\begin{center}
\textbf{\Large Authors' Response to Reviewer 4}
\end{center}

\question{14}{Comment 14 (references) is not adequately addressed. The reference list still contains only 16 entries, essentially unchanged from the original version. The recommended literature was not incorporated, and the coverage remains insufficient for a submission of this scope. The related work should be substantively expanded to reflect the current state of the art, including the previously suggested references.}

\answer
Thank you for pointing out that the recommended studies were missing from the manuscript. Section~3.2 (\emph{Other algorithms}) now cites and discusses all three papers in the text. Each also appears in the reference list.

\begin{enumerate}[leftmargin=2em]
    \item \textbf{Palakonda et al. (2022), \href{https://doi.org/10.1016/j.ins.2022.05.119}{Pi-MOEA}.} The manuscript explains how Pi-MOEA adapts neighborhoods of population members to identify pivot solutions. It then specifies what AGVNS adapts in the routing problem: the order in which VNS explores route moves, using the record of successful mutations.
    \item \textbf{Palakonda et al. (2018), \href{https://doi.org/10.1109/CEC.2018.8477809}{Differential evolution with stochastic selection}.} We added the bibliography entry and discuss how stochastic selection compares candidate solutions under uncertain inputs in smart grid optimization. The revised text identifies the observed number of unallocated orders as the signal used to control crossover in AGVNS at each decision epoch.
    \item \textbf{Tursunboev et al. (2024), \href{https://doi.org/10.1016/j.apenergy.2024.122847}{Evolutionary hybrid deep learning}.} We explain how evolutionary optimization tunes a model and its detection window in that study. The manuscript then specifies how AGVNS links its own stages by passing route plans and mutation records from evolution to VNS.
\end{enumerate}

The revised section also cites the DPDP survey by Cai et al.\ (2023), hierarchical dispatching under LIFO constraints by Du et al.\ (2023), and cost approximation on the ICAPS benchmark by Horv\'ath et al.\ (2025). The adaptive large neighborhood search of Ropke and Pisinger (2006) provides routing context for using past operator performance. The new discussion follows the original review of DPDP methods and leads into the AGVNS methodology.

\question{5}{Comment 5 (runtime reporting) is only partially addressed. The clarification that the 10-minute budget is per-instance, and the added complexity analysis in Section 4.4, are useful, but the requested runtime results are still not reported. Given the fixed-budget dynamic setting, an actual runtime or time-to-solution comparison across methods should be provided rather than a complexity analysis alone.}

% TODO: Add the response to Comment 5 when the runtime revision is ready.

\end{document}
